% ===== 模型假设 =====
\section{模型假设}

为简化问题，本文做出以下假设：

\begin{itemize}[itemindent=2em]
  \item \textbf{假设1 [强]}：……。违反后果：如果该假设不成立，……将……。
  \item \textbf{假设2 [强]}：……。违反后果：……。
  \item \textbf{假设3 [弱]}：……。违反后果：……。
  \item \textbf{假设4 [弱]}：……。违反后果：……。
  \item \textbf{假设5 [弱]}：……。违反后果：……。
\end{itemize}

% 注意：假设编号与问题编号无关。假设是对建模条件的简化声明。
% [强]假设：不成立则模型完全失效（如"忽略非线性效应"在有强非线性时）
% [弱]假设：不成立则模型精度下降但结构仍可用（如"数据服从正态分布"）
